\scanpage{15}
\mspage{13}
qui par passage aux ind-objets redonne le \deleted{foncteur} morphisme précédent, ou ce qui revient au $=$, tel que pour tout $m,n\in\mathbb N$ (de $=$ parité), il existe un $N$ \inserted{$\geq m+n$} assez grand tel que l'on ait un diagramme
\diagram{\begin{tikzcd}[ampersand replacement=\&,column sep=1.9em,scale=.85,transform shape,row sep=1.1em,cells={font=\scriptsize}]
\underline{\mathrm{Ch}}_{\Sigma,m}(X;x,y)\times\underline{\mathrm{Ch}}_{\Sigma,n}(X;y,z)\arrow[r,"{*_{m,n}}"]\arrow[dd,hook]\arrow[dr] \& \underline{\mathrm{Ch}}_{\Sigma,m+n}(X;x,z)\arrow[d,hook]\\
\& \underline{\mathrm{Ch}}_{\Sigma,N}(X;x,z)\arrow[d,hook]\\
\underline{\mathrm{Ch}}_\Sigma(X;x,y)\times\underline{\mathrm{Ch}}_\Sigma(X;y,z)\arrow[r,"?"] \& \underline{\mathrm{Ch}}_\Sigma(X;x,z)
\end{tikzcd}}
et où le triangle supérieur \deleted{est} \uline{commutatif dans $\mathrm{Cathot}$}. (Il faut supposer que pour $N$ \inserted{assez} grand, le composé
\diagram{\begin{tikzcd}[ampersand replacement=\&,row sep=.5em,column sep=1.4em]
{}\arrow[d,hook] \& \\ {}\arrow[r,"?"] \& {}
\end{tikzcd}}
se factorise par $\underline{\mathrm{Ch}}_{\Sigma,N}(X;x,z)$ \ldots)

Pour voir que les foncteurs composition $*_{m,n}$ \deleted{\unsure{forment}} \inserted{définissent} un hom de ind-objets dans $(\mathrm{Cathot})$, il suffit de le prouver dans \unsure{les|le} \deleted{\unsure{deux}} cas \unsure{particuliers|particulier} $(m',n')=(m+1,n+1)$. On va faire la figure en supposant $m,n$ \uline{impairs}.

\deleted{\diagram{\begin{tikzcd}[ampersand replacement=\&,column sep=.17em,scale=.85,transform shape,row sep=.25em,cells={font=\tiny}]
\& \smash{\overbrace{\hspace{2.4em}}^{\text{dég}}} \& \& \deleted{A} \& \& \smash{\overbrace{\hspace{4em}}^{\text{dég.}}} \& \& \& \& \& \& \smash{\overbrace{\hspace{2.7em}}^{\text{dég}}}\\
A_{-n}\arrow[r,"\mathrm{id}"] \& A_{-n} \& A_{-n+1}\arrow[l]\,\cdots \& A_{n-1}\arrow[l]\arrow[r] \& \underset{\deleted{\unsure{B_{-n}}}}{A_n} \& \underset{\substack{\Vert\\B_{-n}}}{A_n}\arrow[l,"\mathrm{id}"']\arrow[r,"\mathrm{id}"] \& B_{-n} \& B_{\unsure{n+1|-n+1}}\arrow[l]\arrow[r] \& \cdots\,B_{\unsure{n-1}} \& B_{\unsure{n}}\arrow[l]\arrow[r] \& B_m \& B_m\arrow[l,"\mathrm{id}"']\\
A_{-n}\arrow[r,"\mathrm{id}"] \& A_{-n} \& A_{-n}\arrow[l,"\mathrm{id}"'] \& \deleted{\unsure{A}}A_{\unsure{n-2|-n-2}}\arrow[r,dash] \& A_{n-1}\\
\& \smash{\underbrace{\hspace{4em}}_{\text{dég}}}
\end{tikzcd}}}

Je m'aperçois \unsure{ici|est} que je me suis \unsure{coulé|couvé} stupidement, la \deleted{\unsure{conjugaison}} \inserted{composition} de deux chemins de longueur \uline{impaire} \inserted{$m,n$} \uline{\unsure{n'est}}

% archive sheet number: 14
